% 7. 準平衡と、回転を失う速さ
\section{準平衡と、回転を失う速さ}\label{sec:equilibrium}

\S\ref{sec:selfint} の率の方程式を解く。自己相互作用が強い($f$ が小さい)と、雲は大きく育つ前に「準平衡」に落ち着き、BH はゆっくりとしか回転を失わない。

\subsection{準平衡の占有数}

$\varepsilon$ の単位の率の方程式(\S\ref{sec:selfint})で、$|322\rangle$ の超放射 $\gamma_{322}$ は非常に小さいので落とす(211 の臨界より上ではこれでよい。臨界より下では $\gamma_{322}$ が主役になる。\S\ref{sec:below211})。スピンの変化への寄与の比 $2\gamma_{322}\varepsilon_{322}/(\gamma_{211}\varepsilon_{211})$ は、GW241011 の主星の例で $4\times10^{-8}$ ほどである(研究ノート E3)。
\begin{align}
\dot\varepsilon_{211}&=\gamma_\mathrm{SR}\varepsilon_{211}-2\gamma_\mathrm{BH}\varepsilon_{211}^2\varepsilon_{322}+\gamma_\infty\varepsilon_{211}\varepsilon_{322}^2,\\
\dot\varepsilon_{322}&=\gamma_\mathrm{BH}\varepsilon_{211}^2\varepsilon_{322}-2\gamma_\infty\varepsilon_{211}\varepsilon_{322}^2.
\end{align}
($\gamma_\mathrm{SR}\equiv\gamma_{211}$ と書いた。)両方を 0 とおく。

\paragraph{手順 1: $|322\rangle$ の式}
$\varepsilon_{211}\varepsilon_{322}\ne0$ で割ると
\begin{equation}
\gamma_\mathrm{BH}\varepsilon_{211}-2\gamma_\infty\varepsilon_{322}=0
\quad\Longrightarrow\quad
\varepsilon_{322}=\frac{\gamma_\mathrm{BH}}{2\gamma_\infty}\,\varepsilon_{211}.
\label{eq:eq322}
\end{equation}

\paragraph{手順 2: $|211\rangle$ の式}
$\varepsilon_{211}$ で割り、式 \eqref{eq:eq322} を代入する。
\begin{align}
0&=\gamma_\mathrm{SR}-2\gamma_\mathrm{BH}\varepsilon_{211}\varepsilon_{322}+\gamma_\infty\varepsilon_{322}^2\notag\\
&=\gamma_\mathrm{SR}-2\gamma_\mathrm{BH}\varepsilon_{211}\cdot\frac{\gamma_\mathrm{BH}\varepsilon_{211}}{2\gamma_\infty}+\gamma_\infty\Bigl(\frac{\gamma_\mathrm{BH}\varepsilon_{211}}{2\gamma_\infty}\Bigr)^2\notag\\
&=\gamma_\mathrm{SR}-\frac{\gamma_\mathrm{BH}^2\varepsilon_{211}^2}{\gamma_\infty}+\frac{\gamma_\mathrm{BH}^2\varepsilon_{211}^2}{4\gamma_\infty}
=\gamma_\mathrm{SR}-\frac34\,\frac{\gamma_\mathrm{BH}^2}{\gamma_\infty}\,\varepsilon_{211}^2.
\end{align}
したがって
\begin{equation}
\boxed{\;\varepsilon_{211}^\mathrm{eq}=\frac{2}{\sqrt3}\,\frac{\sqrt{\gamma_\mathrm{SR}\gamma_\infty}}{\gamma_\mathrm{BH}},\qquad
\varepsilon_{322}^\mathrm{eq}=\frac{1}{\sqrt3}\sqrt{\frac{\gamma_\mathrm{SR}}{\gamma_\infty}}\;}
\label{eq:eq_occ}
\end{equation}
(後者は式 \eqref{eq:eq322} に前者を入れた。)BGLS の式 55a,b と同じである。

\paragraph{$f$ への依存}
$\gamma_\mathrm{SR}$ は $f$ によらず、$\gamma_\infty$、$\gamma_\mathrm{BH}$ はどちらも $(\Mpl/f)^4$ に比例する(式 \ref{eq:gamma_SI})。よって
\begin{equation}
\varepsilon_{211}^\mathrm{eq}\propto\frac{(\Mpl/f)^2}{(\Mpl/f)^4}=\Bigl(\frac{f}{\Mpl}\Bigr)^2,\qquad
\varepsilon_{322}^\mathrm{eq}\propto\Bigl(\frac{f}{\Mpl}\Bigr)^2.
\end{equation}
自己相互作用が強い(小さい $f$)ほど、雲は小さいところで止まる。

\begin{numcheck}
BGLS の例($\alpha=0.01$、$\chi=0.9$、$f=10^{15}\GeV$、主要項の $\gamma_\mathrm{SR}=0.04\,\alpha^8(\chi-2\alpha\rp)\mu$)で確かめる。$\rp=1+\sqrt{0.19}=1.436$、$\chi-2\alpha\rp=0.871$、$(\Mpl/f)^4=(1.221\times10^4)^4=2.22\times10^{16}$。
\begin{align*}
\gamma_\mathrm{SR}/\mu&=0.04\times10^{-16}\times0.871=3.48\times10^{-18},\\
\gamma_\infty/\mu&=1.1\times10^{-8}\times10^{-16}\times2.22\times10^{16}=2.44\times10^{-8},\\
\gamma_\mathrm{BH}/\mu&=4.3\times10^{-7}\times10^{-22}\times2.22\times10^{16}\times1.436=1.37\times10^{-12}.
\end{align*}
式 \eqref{eq:eq_occ} から $\varepsilon_{211}^\mathrm{eq}=1.155\times\sqrt{3.48\times10^{-18}\times2.44\times10^{-8}}/1.37\times10^{-12}=0.246$。BGLS の 0.25 と一致する(研究ノート E3)。
\end{numcheck}

\subsection{準平衡での角運動量とエネルギーの流れ}

単位時間あたり、超放射で $|211\rangle$ に $S=\gamma_\mathrm{SR}\varepsilon_{211}$ の粒子が生まれ、過程 (a)、(b) がそれぞれ $R_\mathrm{BH}$、$R_\infty$ 回起こるとする。準平衡では
\begin{equation}
|211\rangle:\ S-2R_\mathrm{BH}+R_\infty=0,\qquad |322\rangle:\ R_\mathrm{BH}-2R_\infty=0
\quad\Longrightarrow\quad
R_\mathrm{BH}=2R_\infty,\quad S=3R_\infty.
\end{equation}
\begin{itemize}
\item \textbf{角運動量}: BH は単位時間に $S\hbar=3R_\infty\hbar$ を失う。(b) は 1 回で $3\hbar$ を無限遠に運ぶので、運び出す量は $3R_\infty\hbar$。BH から取り出した角運動量は、すべて $m=3$ の波で無限遠に出ていく。雲にはたまらない。
\item \textbf{エネルギー}($E_n\simeq\mu$ の近似): BH は超放射で $S\mu=3R_\infty\mu$ を失い、(a) で $R_\mathrm{BH}(2E_{211}-E_{322})\simeq2R_\infty\mu$ を吸い戻す。差し引き、取り出したエネルギーの 1/3 だけを失う。残りの $R_\infty\mu$ は (b) で無限遠に出ていく。
\end{itemize}

\begin{figure}[H]
\centering
% 準平衡での角運動量とエネルギーの流れ
\begin{tikzpicture}[
  bx/.style={draw, rounded corners=2pt, minimum width=1.5cm, minimum height=0.75cm, font=\small, line width=0.7pt},
  ar/.style={-{Stealth[length=2.2mm]}, line width=0.9pt}]
  \node[bx, fill=ink, text=white] (bh) at (0,0) {BH};
  \node[bx, draw=cA, text=cA] (s1) at (4.0,0) {$|211\rangle$};
  \node[bx, draw=cB, text=cB] (s2) at (8.0,0) {$|322\rangle$};
  \node[bx, draw=cC, text=cC!70!black] (inf) at (12.0,0) {無限遠};
  \draw[ar, ink2] (bh) -- node[above, font=\scriptsize, align=center] {超放射 $S$\\$J$: $3R_\infty\hbar$} (s1);
  \draw[ar, ink2] (s1) -- node[above, font=\scriptsize, align=center] {(a) $R_\mathrm{BH}=2R_\infty$} (s2);
  \draw[ar, ink2] (s2) -- node[above, font=\scriptsize, align=center] {(b) $R_\infty$\\$J$: $3R_\infty\hbar$} (inf);
  \draw[ar, ink2, dashed] (s1.south) to[out=-130, in=-50] node[below, font=\scriptsize, align=center] {(a) で吸い戻すエネルギー $\simeq2R_\infty\mu$(角運動量 0)} (bh.south);
  \draw[ar, ink2, dashed] (s2.north) to[out=140, in=40] node[above, font=\scriptsize] {(b) で $|211\rangle$ に戻る 1 個} (s1.north);
\end{tikzpicture}

\caption{準平衡での角運動量とエネルギーの流れ。BH から $|211\rangle$ に入った角運動量は、$|322\rangle$ を経て $m=3$ の波で外へ出る。エネルギーは 2/3 が BH に戻る。}
\label{fig:flow}
\end{figure}

\subsection{回転を失う速さ}

BH の角運動量を変えるのは超放射だけなので
\begin{equation}
\dot\chi=-\gamma_\mathrm{SR}\,\varepsilon_{211}^\mathrm{eq}=-\frac{2}{\sqrt3}\,\frac{\gamma_\mathrm{SR}^{3/2}\gamma_\infty^{1/2}}{\gamma_\mathrm{BH}}.
\label{eq:chidot}
\end{equation}
$\gamma_\mathrm{SR}\equiv g\,\mu$($g$ は \S\ref{sec:atom} の粒子数の率を $\mu$ で割ったもの)と式 \eqref{eq:gamma_SI} を代入する。
\begin{align}
\gamma_\mathrm{SR}^{3/2}&=g^{3/2}\mu^{3/2},\qquad
\gamma_\infty^{1/2}=\kappa_\infty^{1/2}\alpha^4\Bigl(\frac{\Mpl}{f}\Bigr)^2\mu^{1/2},\notag\\
\frac{\gamma_\mathrm{SR}^{3/2}\gamma_\infty^{1/2}}{\gamma_\mathrm{BH}}
&=\frac{g^{3/2}\mu^{3/2}\cdot\kappa_\infty^{1/2}\alpha^4(\Mpl/f)^2\mu^{1/2}}{\kappa_\mathrm{BH}\alpha^{11}(\Mpl/f)^4\rp\,\mu}
=\mu\Bigl(\frac{f}{\Mpl}\Bigr)^2\frac{g^{3/2}\kappa_\infty^{1/2}}{\kappa_\mathrm{BH}\,\alpha^7\rp}.
\end{align}
まとめると
\begin{equation}
\boxed{\;\dot\chi=\mu\Bigl(\frac{f}{\Mpl}\Bigr)^2H(\alpha,\chi),\qquad
H(\alpha,\chi)=-\frac{2}{\sqrt3}\,\frac{g(\alpha,\chi)^{3/2}\,\kappa_\infty^{1/2}}{\kappa_\mathrm{BH}\,\alpha^7\,\rp}\;}
\label{eq:chidot_H}
\end{equation}
回転を失う速さは $f^2$ に比例する。$f$ が 10 倍小さいと、100 倍遅い。

\begin{numcheck}
上の BGLS の例で $\mu=10^{-12}\eV$($1.519\times10^3\,\mathrm{s}^{-1}$)とすると、$\gamma_\mathrm{SR}=3.48\times10^{-18}\times1519=5.29\times10^{-15}\,\mathrm{s}^{-1}$、$|\dot\chi|=\gamma_\mathrm{SR}\varepsilon_{211}^\mathrm{eq}=1.30\times10^{-15}\,\mathrm{s}^{-1}$。回転を失う時間の目安は $1/|\dot\chi|=7.7\times10^{14}\,\mathrm{s}=2.4\times10^7$ 年。BGLS の式 60 に印刷された $10^7$ 年は、同論文自身の定数から組んだこの値より 2.4 倍小さい(研究ノート E3)。
\end{numcheck}

\subsection{閉じた解}

主要項 $g=\kappa_\mathrm{SR}\,\alpha^8x$($\kappa_\mathrm{SR}\simeq1/24$、$x\equiv\chi-2\alpha\rp$)を使い、$\rp$ と $\alpha$ を一定とみなすと、式 \eqref{eq:chidot_H} は $x$ だけの式になる。$\dd\chi=\dd x$ なので
\begin{equation}
\dot x=-K\,x^{3/2},\qquad
K=\frac{2}{\sqrt3}\,\frac{\kappa_\mathrm{SR}^{3/2}\kappa_\infty^{1/2}}{\kappa_\mathrm{BH}\,\rp}\,\alpha^5\,\mu\Bigl(\frac{f}{\Mpl}\Bigr)^2
\label{eq:xdot}
\end{equation}
($g^{3/2}/\alpha^7=\kappa_\mathrm{SR}^{3/2}\alpha^{12}x^{3/2}/\alpha^7=\kappa_\mathrm{SR}^{3/2}\alpha^5x^{3/2}$)。変数を分けて積分する。
\begin{equation}
\int_{x_0}^{x}x'^{-3/2}\,\dd x'=-K\int_0^t\dd t'
\quad\Longrightarrow\quad
\Bigl[-2x'^{-1/2}\Bigr]_{x_0}^{x}=-Kt
\quad\Longrightarrow\quad
x^{-1/2}=x_0^{-1/2}+\frac{K}{2}t.
\end{equation}
よって
\begin{equation}
\boxed{\;x(t)=\Bigl(x_0^{-1/2}+\frac{K}{2}t\Bigr)^{-2},\qquad
t(x_i\to x_f)=\frac{2}{K}\Bigl(x_f^{-1/2}-x_i^{-1/2}\Bigr)\;}
\label{eq:closedform}
\end{equation}
\begin{itemize}
\item 長い時間では $x\simeq4/(Kt)^2$ で、スピンは臨界値に $t^{-2}$ でしか近づかない。自由場のように臨界値で止まるのではなく、年齢に応じて臨界値の少し上に散らばる。
\item $K\propto\alpha^5$ なので、重い BH ほど速く回転を失う。
\end{itemize}
\begin{caution}
この閉じた解は二つの近似で、数十 \% ずれる。$\rp$ を一定とした($\dd x/\dd\chi=1+2\alpha\chi/\sqrt{1-\chi^2}$ は GW241011 の主星で 1.26)ことと、$g_{11}=1$ とした(実際は 0.6 ほど)ことである。研究のコードは $H(\alpha,\chi)$ をそのまま使っている。閉じた解は、$t^{-2}$ の近づき方と $f^2$ の依存を見るためのものである。
\end{caution}

\begin{figure}[H]
\centering
% 閉じた解 x(t) = (x0^{-1/2} + K t/2)^{-2} によるスピン(説明用)
% chi = 0.338 + x、x0 = 0.442、K(f = 4e12) = 3.4e-8 /yr、K は f^2 に比例
\begin{tikzpicture}
\begin{axis}[note, width=0.66\linewidth, height=0.40\linewidth, xmode=log, xmin=1e4, xmax=1e10, ymin=0.3, ymax=0.82,
  xlabel={年齢 $t$ [yr]}, ylabel={スピン $\chi$}, samples=200, domain=1e4:1e10,
  legend style={font=\scriptsize, at={(0.02,0.03)}, anchor=south west, fill=white, fill opacity=0.85, text opacity=1, draw=none}, legend cell align=left]
  \addplot[cA!45, line width=1pt] {0.338 + (0.442^(-0.5) + 0.5*3.4e-9*x)^(-2)}; \addlegendentry{$f=1.3\times10^{12}\GeV$}
  \addplot[cA, line width=1pt] {0.338 + (0.442^(-0.5) + 0.5*3.4e-8*x)^(-2)}; \addlegendentry{$f=4\times10^{12}\GeV$}
  \addplot[cA!60!black, line width=1pt, dashed] {0.338 + (0.442^(-0.5) + 0.5*3.4e-7*x)^(-2)}; \addlegendentry{$f=1.3\times10^{13}\GeV$}
  \addplot[cB, densely dotted, line width=0.8pt] {0.69};
  \addplot[ink2, densely dashdotted, line width=0.6pt] {0.338};
  \node[font=\scriptsize, color=cB, anchor=south east] at (axis cs:8e9,0.695) {観測の下限 0.69};
  \node[font=\scriptsize, color=ink2, anchor=south east] at (axis cs:8e9,0.343) {$2\alpha\rp$(臨界の近似)};
\end{axis}
\end{tikzpicture}

\caption{閉じた解 \eqref{eq:closedform} によるスピンの時間変化(説明用。$x$ の基準 $2\alpha\rp=0.338$ を一定とした)。GW241011 の主星に近い設定で、$f=4\times10^{12}\GeV$ が $10^7$ 年で 0.69 付近に来るように $K$ を選んだ。$f$ を $\sqrt{10}$ 倍すると曲線は時間方向に 10 倍ずれる。}
\label{fig:spindown}
\end{figure}

\subsection{窓の判定に使う性質}

\paragraph{1. 年齢の比は $f$ によらない}
式 \eqref{eq:chidot_H} では $f$ が前の因子にしかないので、二つの BH が同じ $f$ で回転を失う時間の比
\begin{equation}
R\equiv\frac{t_\mathrm{1G}}{t_\mathrm{2G}}=\frac{\int\dd\chi/|H(\alpha_\mathrm{1G},\chi)|}{\int\dd\chi/|H(\alpha_\mathrm{2G},\chi)|}
\end{equation}
は $f$ によらない。窓があるには、1 世代目と 2 世代目の年齢の比が $R$ 以上でなければならない(研究ノート E3)。

\paragraph{2. $f$ の上限}
ある BH が年齢 $T$ のあいだに許される減少より多く回転を失わないためには、$t(f)\ge T$。$t\propto f^{-2}$ なので、基準の $f_\mathrm{ref}$ での時間 $t_\mathrm{ref}$ から
\begin{equation}
f_\mathrm{max}=f_\mathrm{ref}\sqrt{\frac{t_\mathrm{ref}}{T}}.
\label{eq:fmax}
\end{equation}
GW241011 の主星($0.78\to0.69$、$\mu=0.71\times10^{-12}\eV$)では $f_\mathrm{ref}=10^{13}\GeV$ で $t_\mathrm{ref}=2.0\times10^6$ 年なので、$T=10^7$ 年なら $f_\mathrm{max}=10^{13}\times\sqrt{0.2}=4.5\times10^{12}\GeV$(研究ノート E3)。

\paragraph{3. スピンは三つの量だけで決まる}
$I(\alpha,\chi)\equiv\int_\chi^{\chi_\mathrm{top}}\dd\chi'/|H(\alpha,\chi')|$ とおくと、式 \eqref{eq:chidot_H} を積分して
\begin{equation}
I(\alpha,\chi_T)=I(\alpha,\chi_0)+\tau,\qquad \tau\equiv\mu\Bigl(\frac{f}{\Mpl}\Bigr)^2T.
\end{equation}
年齢 $T$ のあとのスピン $\chi_T$ は $(\alpha,\chi_0,\tau)$ だけで決まる。$I$ を一度表にすれば、集団解析で何百万回でも速く計算できる(研究ノート E10、\texttt{chisr.py})。

\subsection{BH の質量の変化}

準平衡では、BH は取り出した角運動量 $\dd J=-\hbar\,\dd N$ に対し、エネルギーを $\dd M\simeq-\tfrac13\mu\,\dd N$ しか失わない(上の流れ)。$\chi=J/M^2$ の変化は($G=\hbar=1$)
\begin{equation}
\dd\chi=\frac{\dd J}{M^2}-\frac{2J}{M^3}\dd M=-\frac{\dd N}{M^2}+\frac{2\chi}{M}\cdot\frac{\mu\,\dd N}{3}
=-\frac{\dd N}{M^2}\Bigl(1-\frac{2\chi M\mu}{3}\Bigr)=-\frac{\dd N}{M^2}\Bigl(1-\frac{2\chi\alpha}{3}\Bigr).
\end{equation}
質量を固定した式より $(1-2\chi\alpha/3)$ 倍遅い。雲の大きさによらない補正で、GW241011 の主星では 5.4\%(研究ノート E9)。

\subsection{211 の臨界より下: $\gamma_{322}$ を残した固定点}\label{sec:below211}

ここまでの式では、$\gamma_\mathrm{SR}\to0$ でスピンの減少が止まり、BH は 211 の臨界 $\chi_c^{(211)}$ にとどまる。しかし $\gamma_{322}>0$ である限り、BH は 322 の臨界 $\chi_c^{(322)}$ まで回転を失いうる。研究ノート E12 で、この降下が $f\gtrsim10^{14}$--$10^{15}\GeV$ で効くことがわかったので、$\gamma_{322}$ を残して解き直す。$\varepsilon_1\equiv\varepsilon_{211}$、$\varepsilon_2\equiv\varepsilon_{322}$、$\gamma_1\equiv\gamma_{211}$、$\gamma_2\equiv\gamma_{322}$ と書く。
\begin{align}
\dot\varepsilon_1&=\gamma_1\varepsilon_1-2\gamma_\mathrm{BH}\varepsilon_1^2\varepsilon_2+\gamma_\infty\varepsilon_1\varepsilon_2^2,\\
\dot\varepsilon_2&=\gamma_2\varepsilon_2+\gamma_\mathrm{BH}\varepsilon_1^2\varepsilon_2-2\gamma_\infty\varepsilon_1\varepsilon_2^2,\\
\dot\chi&=-\gamma_1\varepsilon_1-2\gamma_2\varepsilon_2.
\end{align}

\paragraph{手順 1: 固定点でのスピンの変化}
$\dot\varepsilon_1=\dot\varepsilon_2=0$ から
\begin{equation}
\gamma_1\varepsilon_1=2\gamma_\mathrm{BH}\varepsilon_1^2\varepsilon_2-\gamma_\infty\varepsilon_1\varepsilon_2^2,\qquad
\gamma_2\varepsilon_2=-\gamma_\mathrm{BH}\varepsilon_1^2\varepsilon_2+2\gamma_\infty\varepsilon_1\varepsilon_2^2 .
\end{equation}
これを $\dot\chi$ に入れると、$\gamma_\mathrm{BH}$ の項が消えて
\begin{equation}
\dot\chi=-(2\gamma_\mathrm{BH}\varepsilon_1^2\varepsilon_2-\gamma_\infty\varepsilon_1\varepsilon_2^2)-2(-\gamma_\mathrm{BH}\varepsilon_1^2\varepsilon_2+2\gamma_\infty\varepsilon_1\varepsilon_2^2)
=-3\gamma_\infty\varepsilon_1\varepsilon_2^2 .
\label{eq:chidot_fp}
\end{equation}
$322\times322\to211\times\infty$ が 1 回起きるたびに、$m=2+2-1=3$ の角運動量が外へ出る(上の流れの図と同じ)。

\paragraph{手順 2: 固定点を解く}
$|322\rangle$ の式を $\varepsilon_2$ で割ると $\gamma_2+\gamma_\mathrm{BH}\varepsilon_1^2-2\gamma_\infty\varepsilon_1\varepsilon_2=0$ なので、$u\equiv\gamma_\mathrm{BH}\varepsilon_1^2$ とおいて
\begin{equation}
\varepsilon_2=\frac{\gamma_2+u}{2\gamma_\infty\varepsilon_1}.
\end{equation}
$|211\rangle$ の式を $\varepsilon_1$ で割った $\gamma_1-2\gamma_\mathrm{BH}\varepsilon_1\varepsilon_2+\gamma_\infty\varepsilon_2^2=0$ に入れる。$\gamma_\mathrm{BH}\varepsilon_1\varepsilon_2=\gamma_\mathrm{BH}(\gamma_2+u)/(2\gamma_\infty)$、$\gamma_\infty\varepsilon_2^2=(\gamma_2+u)^2/(4\gamma_\infty\varepsilon_1^2)=\gamma_\mathrm{BH}(\gamma_2+u)^2/(4\gamma_\infty u)$ なので
\begin{equation}
\gamma_1-\frac{\gamma_\mathrm{BH}(\gamma_2+u)}{\gamma_\infty}+\frac{\gamma_\mathrm{BH}(\gamma_2+u)^2}{4\gamma_\infty u}=0 .
\end{equation}
両辺に $4\gamma_\infty u/\gamma_\mathrm{BH}$ を掛けて整理すると
\begin{equation}
3u^2+\Bigl(2\gamma_2-\frac{4\gamma_\infty\gamma_1}{\gamma_\mathrm{BH}}\Bigr)u-\gamma_2^2=0 .
\label{eq:fp_quad}
\end{equation}
二つの根の積は $-\gamma_2^2/3<0$ なので、正の根はちょうど一つある。

\paragraph{手順 3: 確かめ($\gamma_2\to0$)}
式 \eqref{eq:fp_quad} は $u=4\gamma_\infty\gamma_1/(3\gamma_\mathrm{BH})$ になり、$\varepsilon_1^2=u/\gamma_\mathrm{BH}=4\gamma_1\gamma_\infty/(3\gamma_\mathrm{BH}^2)$。式 \eqref{eq:eq_occ} に戻る。

\paragraph{手順 4: 211 の臨界より下($\gamma_1<0$)}
$\gamma_1=-|\gamma_1|$ とする。臨界のごく近く(下の注意)を除けば $|\gamma_1|\gamma_\infty/\gamma_\mathrm{BH}\gg\gamma_2$ なので、式 \eqref{eq:fp_quad} の小さい正の根は
\begin{equation}
u\simeq\frac{\gamma_2^2\gamma_\mathrm{BH}}{4|\gamma_1|\gamma_\infty},\qquad
\varepsilon_1=\sqrt{\frac{u}{\gamma_\mathrm{BH}}}=\frac{\gamma_2}{2\sqrt{|\gamma_1|\gamma_\infty}},\qquad
\varepsilon_2\simeq\frac{\gamma_2}{2\gamma_\infty\varepsilon_1}=\sqrt{\frac{|\gamma_1|}{\gamma_\infty}} .
\end{equation}
$\varepsilon_2$ は、$322\times322\to211\times\infty$ による $|211\rangle$ の補充($\gamma_\infty\varepsilon_2^2$)が、$|211\rangle$ の吸収 $|\gamma_1|$ とつり合う値である。式 \eqref{eq:chidot_fp} に入れて
\begin{equation}
\dot\chi\simeq-3\gamma_\infty\cdot\frac{\gamma_2}{2\sqrt{|\gamma_1|\gamma_\infty}}\cdot\frac{|\gamma_1|}{\gamma_\infty}=-\frac32\,\gamma_2\sqrt{\frac{|\gamma_1|}{\gamma_\infty}} .
\end{equation}
\begin{itemize}
\item $\gamma_\infty\propto(\Mpl/f)^4$ なので $\sqrt{1/\gamma_\infty}\propto(f/\Mpl)^2$。この速さも $\mu(f/\Mpl)^2$ に比例し、$\gamma_\mathrm{BH}$ によらない。
\item そのため、前の小節の $I(\alpha,\chi)$ の方法はそのまま使え、表を $\chi_c^{(322)}$ まで延ばせばよい。$\gamma_2\to0$ の $\chi_c^{(322)}$ で $\dot\chi\to0$ となり、BH はそこで止まる。
\end{itemize}

\begin{caution}
回転を失う速さが最も遅いのは、211 の臨界のすぐ上と下、$|\gamma_1|\sim\gamma_2\gamma_\mathrm{BH}/\gamma_\infty$ の狭い範囲(スピンにして $10^{-8}$ ほど)である。数値で積分するときは、ここを細かく刻まないと時間を短く見積もる。
\end{caution}

\begin{numcheck}
GW241011 の主星($\mu=0.71\times10^{-12}\eV$、$\chi_0=0.78$、$\chi_c^{(211)}=0.399$、$\chi_c^{(322)}=0.206$)で、$f=10^{16}\GeV$ では $10^9$ 年後のスピンは 0.220。運動論の ODE は 0.219 である。$f\le10^{13}\GeV$ では $10^{10}$ 年たっても 211 の臨界より下に行かない(研究ノート E12)。
\end{numcheck}

\subsection{いつ準平衡が成り立つか}

\begin{itemize}
\item \textbf{準平衡に早く着くこと}: 雲は 1 個の量子から自由に育ち、$|322\rangle$ が増えると止まる。止まる時点の $\varepsilon_{211}$ は自由場の飽和量よりずっと小さく、それまでに失う回転は無視できる。GW241011 の主星($f=4\times10^{12}\GeV$)では約 30 年で準平衡の 1\% 以内に入る(研究ノート E5 の ODE)。
\item \textbf{二準位で足りること}: $\alpha\gtrsim0.15$ では $|433\rangle$ なども育ち、二準位の近似が崩れる(\arxiv{2412.03655})。
\item \textbf{bosenova に届かないこと}: 引力が強すぎると雲がつぶれる(bosenova)。準平衡の雲はその閾値の 5\% ほどで届かない(研究ノート E3)。
\end{itemize}

\begin{summary}
\begin{itemize}
\item 準平衡の占有数 $\varepsilon_{211}^\mathrm{eq}=(2/\sqrt3)\sqrt{\gamma_\mathrm{SR}\gamma_\infty}/\gamma_\mathrm{BH}\propto f^2$。
\item 取り出した角運動量はすべて $m=3$ の波で外へ出る。エネルギーは 2/3 が BH に戻る。
\item 回転を失う速さ $\dot\chi=\mu(f/\Mpl)^2H(\alpha,\chi)\propto f^2$。主要項では $x(t)=(x_0^{-1/2}+Kt/2)^{-2}$。
\item 年齢の比は $f$ によらず、$f$ の上限は $f_\mathrm{max}=f_\mathrm{ref}\sqrt{t_\mathrm{ref}/T}$、スピンは $(\alpha,\chi_0,\tau)$ で決まる。
\item $\gamma_{322}$ を残すと、211 の臨界より下でも $\dot\chi\simeq-\tfrac32\gamma_{322}\sqrt{|\gamma_{211}|/\gamma_\infty}\propto f^2$ で 322 の臨界まで下がる。
\end{itemize}
\end{summary}

\begin{tocode}
コード: \texttt{si.eps\_eq} が式 \eqref{eq:eq_occ}、\texttt{si.chidot} が式 \eqref{eq:chidot}(\texttt{MASS\_CORR} で質量の補正)、\texttt{si.t\_spindown} が時間の積分、\texttt{chisr.H} が式 \eqref{eq:fp_quad}・\eqref{eq:chidot_fp} の固定点(\S\ref{sec:below211})、\texttt{chisr.chi\_after} が $I$ の表。
\end{tocode}
